\documentclass{article}
\usepackage{amsmath, amssymb, amsthm}
\usepackage{hyperref}
\usepackage{geometry}
\geometry{margin=1in}

\title{Erd\H{o}s Problem \#456}
\date{Last updated: 2025-08-31}
\author{Source: erdosproblems.com}

\begin{document}
\maketitle

\noindent\textbf{Status:} OPEN \quad \textbf{No prize}

\noindent\textbf{Tags:} number theory

\medskip
\noindent\textbf{Problem Statement:}

Let $p_n$ be the smallest prime $\equiv 1\pmod{n}$ and let $m_n$ be the smallest integer such that $n\mid \phi(m_n)$.

Is it true that $m_n<p_n$ for almost all $n$? Does $p_n/m_n\to \infty$ for almost all $n$? Are there infinitely many primes $p$ such that $p-1$ is the only $n$ for which $m_n=p$?

\bigskip
\noindent\textbf{Remarks:}

Linnik's theorem implies that $p_n\leq n^{O(1)}$. It is trivial that $m_n\leq p_n$ always.

If $n=q-1$ for some prime $q$ then $m_n=p_n$. Erd\H{o}s \cite{Er79e} writes it is 'easy to show' that for infinitely many $n$ we have $m_n <p_n$, and that $m_n/n\to \infty$ for almost all $n$.

van Doorn in the comments has noted that if $n=2^{2k+1}$ then $m_n\leq 2n$ and $p_n\geq 2n+1$.

\bigskip
\noindent\textbf{References:}

[Er79e] Erd\H{o}s, Paul, Some unconventional problems in number theory. Ast\'{e}risque (1979), 73-82.

\bigskip
\noindent\small{Source: \url{https://www.erdosproblems.com/456}}
\end{document}
